\section{Experiments}
\label{sec:results}

\subsection{Dataset and Protocol}
All results in this draft follow a unified \textbf{train/validation-only} protocol:
\begin{itemize}
  \item target split mode: \texttt{table\_group};
  \item split ratio: \textbf{1:5} (train:val);
  \item no test split (test=0).
\end{itemize}

The dataset used here is \texttt{output/dataset\_table\_group\_train\_val\_no\_test.pkl}
with target-domain counts:
train=1274, val=6270, test=0.

\subsection{Core Multi-Seed Comparison}
We report multi-seed statistics from
\texttt{paper\_notest\_core\_summary.csv}.

\begin{table}[htbp]
\centering
\small
\setlength{\tabcolsep}{4pt}
\begin{tabular}{lcccc}
\toprule
Method & NSeed & Train $R^2$ (mean+/-std) & Val $R^2$ (mean+/-std) & Val Loss (mean+/-std) \\
\midrule
Baseline (shared linear) & 5 & 0.9024+/-0.0086 & 0.8726+/-0.0067 & 0.121635+/-0.006428 \\
TCDR (16,96) & 5 & 0.9586+/-0.0085 & 0.9065+/-0.0168 & 0.089249+/-0.016086 \\
TCDR (32,64) & 5 & \textbf{0.9582+/-0.0216} & \textbf{0.9234+/-0.0175} & \textbf{0.073105+/-0.016755} \\
\bottomrule
\end{tabular}
\caption{Core architecture comparison under 1:5 train/val protocol (no test split).}
\label{tab:notest_core}
\end{table}

Both TCDR variants significantly outperform the baseline on validation metrics.
In the current no-test setting, TCDR (32,64) gives the strongest mean validation score.

\subsection{One-Factor Ablation (Main: T32x64)}
With TCDR (32,64) as the fixed backbone, we vary one factor at a time and report
multi-seed statistics from \texttt{paper\_t32\_onefactor\_summary.csv}.

\begin{table}[htbp]
\centering
\small
\setlength{\tabcolsep}{4pt}
\begin{tabular}{lcccc}
\toprule
Method & NSeed & Train $R^2$ (mean+/-std) & Val $R^2$ (mean+/-std) & Val Loss (mean+/-std) \\
\midrule
T32x64 + mode\_base(enrich\_on) & 5 & 0.9682+/-0.0087 & \textbf{0.9298+/-0.0073} & \textbf{0.066996+/-0.006930} \\
T32x64 + enrich\_off(auto) & 5 & 0.9682+/-0.0087 & \textbf{0.9298+/-0.0073} & \textbf{0.066996+/-0.006930} \\
T32x64 + parasitic\_replace & 5 & 0.9694+/-0.0080 & 0.9294+/-0.0084 & 0.067451+/-0.008056 \\
T32x64 + dual\_readout(mean) & 5 & 0.9582+/-0.0216 & 0.9234+/-0.0175 & 0.073105+/-0.016755 \\
T32x64 + mode\_base+topology\_off & 5 & 0.9006+/-0.0068 & 0.8720+/-0.0034 & 0.122180+/-0.003232 \\
\bottomrule
\end{tabular}
\caption{One-factor ablations on T32x64 under the same 1:5 no-test protocol.}
\label{tab:notest_component}
\end{table}

\subsection{Current Takeaway}
\begin{itemize}
  \item TCDR remains the primary gain source in no-test OOD validation.
  \item TCDR (32,64) is currently the strongest candidate under this protocol.
  \item In T32x64 ablation, the strongest settings are mode\_base(enrich\_on) and enrich\_off(auto), with identical statistics.
  \item Turning topology expert off causes a large drop (Val $R^2$: 0.9298 $\rightarrow$ 0.8720), confirming topology-conditioned residual as a key contributor.
  \item T16x96 one-factor results are kept as supplementary evidence (see \texttt{paper\_notest\_component\_summary.csv}).
\end{itemize}
